\subsection{Baseline Methods for Performance Evaluation}

To benchmark the effectiveness of the proposed method, the antenna orientation scheme's performance will be compared with that of other antenna alignment strategies, known as baselines, where each one represents a different approach in terms of adaptiveness, computational complexity, and convergence properties.

\subsubsection{Exhaustive Angular Scan Approach}

With exhaustive angular scan, the antenna will go through a set of predetermined angles by moving from one orientation to the next while performing RSSI measurements at a fixed period of time for each orientation. Once a whole scan has been performed, the antenna will face the optimal orientation defined by the maximum RSSI value found during the procedure.

Even though this method is guaranteed to find the optimal angle within the resolution of the scan, its high computational and mechanical cost can lead to long antenna orientation times.

\subsubsection{Hill-Climbing Based on RSSI Values}

The hill-climbing approach based on RSSI values uses a simple local greedy strategy to search for the optimal orientation by applying incremental changes in both axis directions, keeping the new orientation if there is an improvement on RSSI or reverting back to previous orientation otherwise.

Although it is less computationally expensive than the exhaustive scan approach, this method suffers from local optima convergence problems.